% MCQ Questions for Chapter 05
\section{বহুনির্বাচনি প্রশ্ন (MCQ)}

\subsection*{\textbf{\underline{প্রোগ্রামিং ভাষার ইতিহাস}}}

\begin{multicols}{2}
\begin{enumerate}

\item প্রথম প্রোগ্রামিং ভাষা কোনটি? \textbf{[ঢাকা-২০২৩]}
\begin{enumerate}[label=\alph*)]
\item Machine Language
\item Assembly Language
\item FORTRAN
\item COBOL
\end{enumerate}

\item Machine Language কোন ধরনের ভাষা? \textbf{[চট্টগ্রাম-২০২৩]}
\begin{enumerate}[label=\alph*)]
\item Low Level Language
\item High Level Language
\item Middle Level Language
\item Natural Language
\end{enumerate}

\item Assembly Language এ কী ব্যবহার করা হয়? \textbf{[রাজশাহী-২০২২]}
\begin{enumerate}[label=\alph*)]
\item Mnemonic Code
\item Binary Code
\item Hexadecimal Code
\item Decimal Code
\end{enumerate}

\item Assembler কী? \textbf{[সিলেট-২০২৩]}
\begin{enumerate}[label=\alph*)]
\item Assembly কে Machine এ রূপান্তরকারী
\item High Level কে Machine এ রূপান্তরকারী
\item Machine কে Assembly তে রূপান্তরকারী
\item Binary কে Decimal এ রূপান্তরকারী
\end{enumerate}

\item Compiler কী? \textbf{[বরিশাল-২০২২]}
\begin{enumerate}[label=\alph*)]
\item পুরো প্রোগ্রাম একসাথে অনুবাদ করে
\item এক লাইন করে অনুবাদ করে
\item শুধু error খুঁজে
\item প্রোগ্রাম লেখে
\end{enumerate}

\item Interpreter কী? \textbf{[দিনাজপুর-২০২৩]}
\begin{enumerate}[label=\alph*)]
\item এক লাইন করে অনুবাদ করে
\item পুরো প্রোগ্রাম একসাথে অনুবাদ করে
\item শুধু error খুঁজে
\item প্রোগ্রাম লেখে
\end{enumerate}

\item FORTRAN এর পূর্ণরূপ কী? \textbf{[কুমিল্লা-২০২২]}
\begin{enumerate}[label=\alph*)]
\item Formula Translation
\item Forward Translation
\item Format Translation
\item Function Translation
\end{enumerate}

\item C ভাষা কে তৈরি করেন? \textbf{[যশোর-২০২৩]}
\begin{enumerate}[label=\alph*)]
\item Dennis Ritchie
\item Bjarne Stroustrup
\item James Gosling
\item Guido van Rossum
\end{enumerate}

\end{enumerate}
\end{multicols}

\subsection*{\textbf{\underline{প্রোগ্রামিং এর মূলনীতি}}}

\begin{multicols}{2}
\begin{enumerate}[resume]

\item Algorithm কী? \textbf{[ঢাকা-২০২২]}
\begin{enumerate}[label=\alph*)]
\item সমস্যা সমাধানের ধাপসমূহ
\item একটি প্রোগ্রামিং ভাষা
\item একটি সফটওয়্যার
\item একটি হার্ডওয়্যার
\end{enumerate}

\item Flowchart কী? \textbf{[চট্টগ্রাম-২০২২]}
\begin{enumerate}[label=\alph*)]
\item Algorithm এর চিত্ররূপ
\item একটি প্রোগ্রাম
\item একটি ডায়াগ্রাম
\item একটি চার্ট
\end{enumerate}

\item Flowchart এ Oval চিহ্ন কী বোঝায়? \textbf{[রাজশাহী-২০২৩]}
\begin{enumerate}[label=\alph*)]
\item Start/Stop
\item Process
\item Decision
\item Input/Output
\end{enumerate}

\item Flowchart এ Rectangle চিহ্ন কী বোঝায়? \textbf{[সিলেট-২০২২]}
\begin{enumerate}[label=\alph*)]
\item Process
\item Start/Stop
\item Decision
\item Input/Output
\end{enumerate}

\item Flowchart এ Diamond চিহ্ন কী বোঝায়? \textbf{[বরিশাল-২০২৩]}
\begin{enumerate}[label=\alph*)]
\item Decision
\item Process
\item Start/Stop
\item Input/Output
\end{enumerate}

\item Flowchart এ Parallelogram চিহ্ন কী বোঝায়? \textbf{[দিনাজপুর-২০২২]}
\begin{enumerate}[label=\alph*)]
\item Input/Output
\item Process
\item Decision
\item Start/Stop
\end{enumerate}

\item Pseudocode কী? \textbf{[কুমিল্লা-২০২৩]}
\begin{enumerate}[label=\alph*)]
\item সাধারণ ভাষায় লেখা প্রোগ্রাম
\item একটি প্রোগ্রামিং ভাষা
\item একটি সফটওয়্যার
\item একটি হার্ডওয়্যার
\end{enumerate}

\item Debugging কী? \textbf{[যশোর-২০২২]}
\begin{enumerate}[label=\alph*)]
\item Error খুঁজে বের করা ও সংশোধন
\item প্রোগ্রাম লেখা
\item প্রোগ্রাম চালানো
\item প্রোগ্রাম সংরক্ষণ
\end{enumerate}

\end{enumerate}
\end{multicols}

\subsection*{\textbf{\underline{C প্রোগ্রামিং বেসিক}}}

\begin{multicols}{2}
\begin{enumerate}[resume]

\item C প্রোগ্রামের মূল ফাংশন কোনটি? \textbf{[ঢাকা-২০২৩]}
\begin{enumerate}[label=\alph*)]
\item main()
\item start()
\item begin()
\item run()
\end{enumerate}

\item C প্রোগ্রামে Header File include করা হয় কীভাবে? \textbf{[চট্টগ্রাম-২০২৩]}
\begin{enumerate}[label=\alph*)]
\item \#include <stdio.h>
\item import stdio.h
\item using stdio.h
\item require stdio.h
\end{enumerate}

\item printf() ফাংশনের কাজ কী? \textbf{[রাজশাহী-২০২২]}
\begin{enumerate}[label=\alph*)]
\item Output প্রদর্শন
\item Input নেওয়া
\item Calculation করা
\item File খোলা
\end{enumerate}

\item scanf() ফাংশনের কাজ কী? \textbf{[সিলেট-২০২৩]}
\begin{enumerate}[label=\alph*)]
\item Input নেওয়া
\item Output প্রদর্শন
\item Calculation করা
\item File খোলা
\end{enumerate}

\item Integer variable declare করা হয় কীভাবে? \textbf{[বরিশাল-২০২২]}
\begin{enumerate}[label=\alph*)]
\item int a;
\item integer a;
\item number a;
\item var a;
\end{enumerate}

\item Float variable declare করা হয় কীভাবে? \textbf{[দিনাজপুর-২০২৩]}
\begin{enumerate}[label=\alph*)]
\item float a;
\item decimal a;
\item real a;
\item double a;
\end{enumerate}

\item Character variable declare করা হয় কীভাবে? \textbf{[কুমিল্লা-২০২২]}
\begin{enumerate}[label=\alph*)]
\item char a;
\item character a;
\item string a;
\item text a;
\end{enumerate}

\item C তে comment লেখা হয় কীভাবে? \textbf{[যশোর-২০২৩]}
\begin{enumerate}[label=\alph*)]
\item // comment বা /* comment */
\item \# comment
\item -- comment
\item ' comment
\end{enumerate}

\end{enumerate}
\end{multicols}

\subsection*{\textbf{\underline{Operators ও Expressions}}}

\begin{multicols}{2}
\begin{enumerate}[resume]

\item Arithmetic Operator কোনটি? \textbf{[ঢাকা-২০২২]}
\begin{enumerate}[label=\alph*)]
\item +, -, *, /, \%
\item ==, !=, <, >
\item \&\&, ||, !
\item =, +=, -=
\end{enumerate}

\item Relational Operator কোনটি? \textbf{[চট্টগ্রাম-২০২২]}
\begin{enumerate}[label=\alph*)]
\item ==, !=, <, >, <=, >=
\item +, -, *, /
\item \&\&, ||, !
\item =, +=, -=
\end{enumerate}

\item Logical Operator কোনটি? \textbf{[রাজশাহী-২০২৩]}
\begin{enumerate}[label=\alph*)]
\item \&\&, ||, !
\item +, -, *, /
\item ==, !=, <, >
\item =, +=, -=
\end{enumerate}

\item Assignment Operator কোনটি? \textbf{[সিলেট-২০২২]}
\begin{enumerate}[label=\alph*)]
\item =, +=, -=, *=, /=
\item +, -, *, /
\item ==, !=, <, >
\item \&\&, ||, !
\end{enumerate}

\item Increment Operator কোনটি? \textbf{[বরিশাল-২০২৩]}
\begin{enumerate}[label=\alph*)]
\item ++
\item --
\item +=
\item *=
\end{enumerate}

\item Decrement Operator কোনটি? \textbf{[দিনাজপুর-২০২২]}
\begin{enumerate}[label=\alph*)]
\item --
\item ++
\item -=
\item /=
\end{enumerate}

\item Modulus Operator কোনটি? \textbf{[কুমিল্লা-২০২৩]}
\begin{enumerate}[label=\alph*)]
\item \%
\item /
\item *
\item //
\end{enumerate}

\item a++ এবং ++a এর পার্থক্য কী? \textbf{[যশোর-২০২২]}
\begin{enumerate}[label=\alph*)]
\item a++ post-increment, ++a pre-increment
\item কোনো পার্থক্য নেই
\item a++ দ্রুত
\item ++a দ্রুত
\end{enumerate}

\end{enumerate}
\end{multicols}

\subsection*{\textbf{\underline{Control Structures}}}

\begin{multicols}{2}
\begin{enumerate}[resume]

\item if statement এর syntax কোনটি? \textbf{[ঢাকা-২০২৩]}
\begin{enumerate}[label=\alph*)]
\item if (condition) \{ \}
\item if condition then
\item if [condition]
\item if condition:
\end{enumerate}

\item if-else statement কখন ব্যবহৃত হয়? \textbf{[চট্টগ্রাম-২০২৩]}
\begin{enumerate}[label=\alph*)]
\item দুটি option এর জন্য
\item একটি option এর জন্য
\item তিনটি option এর জন্য
\item অসীম option এর জন্য
\end{enumerate}

\item switch statement কখন ব্যবহৃত হয়? \textbf{[রাজশাহী-২০২২]}
\begin{enumerate}[label=\alph*)]
\item একাধিক option এর জন্য
\item দুটি option এর জন্য
\item একটি option এর জন্য
\item কোনো option নেই
\end{enumerate}

\item for loop এর syntax কোনটি? \textbf{[সিলেট-২০২৩]}
\begin{enumerate}[label=\alph*)]
\item for (init; condition; increment)
\item for condition
\item for [init, condition, increment]
\item for init to condition
\end{enumerate}

\item while loop কখন ব্যবহৃত হয়? \textbf{[বরিশাল-২০২২]}
\begin{enumerate}[label=\alph*)]
\item iteration সংখ্যা অজানা থাকলে
\item iteration সংখ্যা জানা থাকলে
\item শুধু একবার চালাতে
\item কখনো চালাতে হবে না
\end{enumerate}

\item do-while loop এর বৈশিষ্ট্য কী? \textbf{[দিনাজপুর-২০২৩]}
\begin{enumerate}[label=\alph*)]
\item কমপক্ষে একবার চলবে
\item কখনো নাও চলতে পারে
\item অসীমবার চলবে
\item কখনো চলবে না
\end{enumerate}

\item break statement এর কাজ কী? \textbf{[কুমিল্লা-২০২২]}
\begin{enumerate}[label=\alph*)]
\item loop থেকে বের হওয়া
\item loop চালু করা
\item loop বন্ধ করা
\item loop পুনরায় শুরু করা
\end{enumerate}

\item continue statement এর কাজ কী? \textbf{[যশোর-২০২৩]}
\begin{enumerate}[label=\alph*)]
\item বর্তমান iteration skip করা
\item loop থেকে বের হওয়া
\item loop বন্ধ করা
\item loop শুরু করা
\end{enumerate}

\end{enumerate}
\end{multicols}

\subsection*{\textbf{\underline{Array ও Function}}}

\begin{multicols}{2}
\begin{enumerate}[resume]

\item Array কী? \textbf{[ঢাকা-২০২২]}
\begin{enumerate}[label=\alph*)]
\item একই ধরনের ডেটার সংগ্রহ
\item ভিন্ন ধরনের ডেটার সংগ্রহ
\item একটি ফাংশন
\item একটি অপারেটর
\end{enumerate}

\item Array declare করা হয় কীভাবে? \textbf{[চট্টগ্রাম-২০২২]}
\begin{enumerate}[label=\alph*)]
\item int arr[10];
\item array arr[10];
\item int arr(10);
\item array(10) arr;
\end{enumerate}

\item Array এর index কত থেকে শুরু হয়? \textbf{[রাজশাহী-২০২৩]}
\begin{enumerate}[label=\alph*)]
\item ০
\item ১
\item -১
\item যেকোনো সংখ্যা
\end{enumerate}

\item 2D Array declare করা হয় কীভাবে? \textbf{[সিলেট-২০২২]}
\begin{enumerate}[label=\alph*)]
\item int arr[3][4];
\item int arr[3,4];
\item int arr(3)(4);
\item array arr[3][4];
\end{enumerate}

\item Function কী? \textbf{[বরিশাল-২০২৩]}
\begin{enumerate}[label=\alph*)]
\item পুনর্ব্যবহারযোগ্য code block
\item একটি variable
\item একটি operator
\item একটি loop
\end{enumerate}

\item Function declare করা হয় কীভাবে? \textbf{[দিনাজপুর-২০২২]}
\begin{enumerate}[label=\alph*)]
\item return\_type function\_name(parameters)
\item function function\_name
\item def function\_name
\item func function\_name
\end{enumerate}

\item void function কী return করে? \textbf{[কুমিল্লা-২০২৩]}
\begin{enumerate}[label=\alph*)]
\item কিছুই না
\item ০
\item ১
\item NULL
\end{enumerate}

\item Recursion কী? \textbf{[যশোর-২০২২]}
\begin{enumerate}[label=\alph*)]
\item ফাংশন নিজেকে call করা
\item একটি loop
\item একটি array
\item একটি variable
\end{enumerate}

\end{enumerate}
\end{multicols}

\subsection*{\textbf{\underline{উত্তরমালা}}}

\begin{multicols}{4}
\begin{enumerate}
\item a
\item a
\item a
\item a
\item a
\item a
\item a
\item a
\item a
\item a
\item a
\item a
\item a
\item a
\item a
\item a
\item a
\item a
\item a
\item a
\item a
\item a
\item a
\item a
\item a
\item a
\item a
\item a
\item a
\item a
\item a
\item a
\item a
\item a
\item a
\item a
\item a
\item a
\item a
\item a
\item a
\item a
\item a
\item a
\item a
\item a
\item a
\item a
\item a
\item a
\item a
\item a
\item a
\item a
\item a
\item a
\end{enumerate}
\end{multicols}

\subsection*{\textbf{\underline{Pointer ও String}}}

\begin{multicols}{2}
\begin{enumerate}[resume]

\item Pointer কী? \textbf{[ঢাকা-২০২৩]}
\begin{enumerate}[label=\alph*)]
\item মেমরি address সংরক্ষণকারী variable
\item একটি array
\item একটি function
\item একটি operator
\end{enumerate}

\item Pointer declare করা হয় কীভাবে? \textbf{[চট্টগ্রাম-২০২২]}
\begin{enumerate}[label=\alph*)]
\item int *ptr;
\item int ptr*;
\item pointer ptr;
\item int \&ptr;
\end{enumerate}

\item Address operator কোনটি? \textbf{[সিলেট-২০২৩]}
\begin{enumerate}[label=\alph*)]
\item \&
\item *
\item \#
\item @
\end{enumerate}

\item Dereference operator কোনটি? \textbf{[রাজশাহী-২০২২]}
\begin{enumerate}[label=\alph*)]
\item *
\item \&
\item \#
\item @
\end{enumerate}

\item NULL pointer কী? \textbf{[বরিশাল-২০২৩]}
\begin{enumerate}[label=\alph*)]
\item যে pointer কোনো address point করে না
\item যে pointer সব address point করে
\item একটি array
\item একটি function
\end{enumerate}

\item String কী? \textbf{[দিনাজপুর-২০২২]}
\begin{enumerate}[label=\alph*)]
\item Character এর array
\item Integer এর array
\item Float এর array
\item একটি function
\end{enumerate}

\item String declare করা হয় কীভাবে? \textbf{[কুমিল্লা-২০২৩]}
\begin{enumerate}[label=\alph*)]
\item char str[20];
\item string str;
\item text str;
\item str[20];
\end{enumerate}

\item String এর শেষে কী থাকে? \textbf{[যশোর-২০২২]}
\begin{enumerate}[label=\alph*)]
\item \textbackslash 0 (null character)
\item \textbackslash n (newline)
\item \textbackslash t (tab)
\item space
\end{enumerate}

\item strlen() function কী করে? \textbf{[ঢাকা-২০২৩]}
\begin{enumerate}[label=\alph*)]
\item String এর দৈর্ঘ্য বের করে
\item String copy করে
\item String তুলনা করে
\item String যোগ করে
\end{enumerate}

\item strcpy() function কী করে? \textbf{[চট্টগ্রাম-২০২২]}
\begin{enumerate}[label=\alph*)]
\item String copy করে
\item String এর দৈর্ঘ্য বের করে
\item String তুলনা করে
\item String যোগ করে
\end{enumerate}

\item strcmp() function কী করে? \textbf{[সিলেট-২০২৩]}
\begin{enumerate}[label=\alph*)]
\item String তুলনা করে
\item String copy করে
\item String এর দৈর্ঘ্য বের করে
\item String যোগ করে
\end{enumerate}

\item strcat() function কী করে? \textbf{[রাজশাহী-২০২২]}
\begin{enumerate}[label=\alph*)]
\item String যোগ করে
\item String copy করে
\item String তুলনা করে
\item String এর দৈর্ঘ্য বের করে
\end{enumerate}

\end{enumerate}
\end{multicols}

\subsection*{\textbf{\underline{Structure ও File}}}

\begin{multicols}{2}
\begin{enumerate}[resume]

\item Structure কী? \textbf{[বরিশাল-২০২৩]}
\begin{enumerate}[label=\alph*)]
\item ভিন্ন ধরনের ডেটার সংগ্রহ
\item একই ধরনের ডেটার সংগ্রহ
\item একটি function
\item একটি operator
\end{enumerate}

\item Structure declare করা হয় কীভাবে? \textbf{[দিনাজপুর-২০২২]}
\begin{enumerate}[label=\alph*)]
\item struct name \{ \};
\item structure name \{ \};
\item class name \{ \};
\item object name \{ \};
\end{enumerate}

\item Structure variable declare করা হয় কীভাবে? \textbf{[কুমিল্লা-২০২৩]}
\begin{enumerate}[label=\alph*)]
\item struct name var;
\item name var;
\item var struct name;
\item structure var;
\end{enumerate}

\item Structure member access করা হয় কীভাবে? \textbf{[যশোর-২০২২]}
\begin{enumerate}[label=\alph*)]
\item var.member
\item var->member
\item var[member]
\item var(member)
\end{enumerate}

\item File কী? \textbf{[ঢাকা-২০২৩]}
\begin{enumerate}[label=\alph*)]
\item ডেটা সংরক্ষণের মাধ্যম
\item একটি variable
\item একটি function
\item একটি operator
\end{enumerate}

\item File open করার function কোনটি? \textbf{[চট্টগ্রাম-২০২২]}
\begin{enumerate}[label=\alph*)]
\item fopen()
\item open()
\item file\_open()
\item start()
\end{enumerate}

\item File close করার function কোনটি? \textbf{[সিলেট-২০২৩]}
\begin{enumerate}[label=\alph*)]
\item fclose()
\item close()
\item file\_close()
\item end()
\end{enumerate}

\item File থেকে পড়ার function কোনটি? \textbf{[রাজশাহী-২০২২]}
\begin{enumerate}[label=\alph*)]
\item fscanf(), fgets()
\item read()
\item input()
\item get()
\end{enumerate}

\item File এ লেখার function কোনটি? \textbf{[বরিশাল-২০২৩]}
\begin{enumerate}[label=\alph*)]
\item fprintf(), fputs()
\item write()
\item output()
\item put()
\end{enumerate}

\item File mode "r" কী বোঝায়? \textbf{[দিনাজপুর-২০২২]}
\begin{enumerate}[label=\alph*)]
\item Read mode
\item Write mode
\item Append mode
\item Binary mode
\end{enumerate}

\item File mode "w" কী বোঝায়? \textbf{[কুমিল্লা-২০২৩]}
\begin{enumerate}[label=\alph*)]
\item Write mode
\item Read mode
\item Append mode
\item Binary mode
\end{enumerate}

\item File mode "a" কী বোঝায়? \textbf{[যশোর-২০২২]}
\begin{enumerate}[label=\alph*)]
\item Append mode
\item Read mode
\item Write mode
\item Binary mode
\end{enumerate}

\end{enumerate}
\end{multicols}

\subsection*{\textbf{\underline{প্রোগ্রামিং কনসেপ্ট}}}

\begin{multicols}{2}
\begin{enumerate}[resume]

\item Variable কী? \textbf{[ঢাকা-২০২৩]}
\begin{enumerate}[label=\alph*)]
\item ডেটা সংরক্ষণের স্থান
\item একটি function
\item একটি operator
\item একটি loop
\end{enumerate}

\item Constant কী? \textbf{[চট্টগ্রাম-২০২২]}
\begin{enumerate}[label=\alph*)]
\item যার মান পরিবর্তন হয় না
\item যার মান পরিবর্তন হয়
\item একটি function
\item একটি operator
\end{enumerate}

\item Data Type কী? \textbf{[সিলেট-২০২৩]}
\begin{enumerate}[label=\alph*)]
\item ডেটার ধরন নির্ধারণ
\item ডেটা সংরক্ষণ
\item ডেটা প্রসেসিং
\item ডেটা মুছে ফেলা
\end{enumerate}

\item Scope কী? \textbf{[রাজশাহী-২০২২]}
\begin{enumerate}[label=\alph*)]
\item Variable এর ব্যবহারযোগ্য এলাকা
\item Variable এর মান
\item Variable এর নাম
\item Variable এর ধরন
\end{enumerate}

\item Local Variable কী? \textbf{[বরিশাল-২০২৩]}
\begin{enumerate}[label=\alph*)]
\item Function এর ভিতরে declare করা
\item Function এর বাইরে declare করা
\item সব জায়গায় ব্যবহারযোগ্য
\item কোথাও ব্যবহারযোগ্য নয়
\end{enumerate}

\item Global Variable কী? \textbf{[দিনাজপুর-২০২২]}
\begin{enumerate}[label=\alph*)]
\item Function এর বাইরে declare করা
\item Function এর ভিতরে declare করা
\item শুধু main() এ ব্যবহারযোগ্য
\item কোথাও ব্যবহারযোগ্য নয়
\end{enumerate}

\item Parameter কী? \textbf{[কুমিল্লা-২০২৩]}
\begin{enumerate}[label=\alph*)]
\item Function এ পাঠানো মান
\item Function এর return মান
\item একটি variable
\item একটি operator
\end{enumerate}

\item Argument কী? \textbf{[যশোর-২০২২]}
\begin{enumerate}[label=\alph*)]
\item Function call এ পাঠানো মান
\item Function এর return মান
\item একটি variable
\item একটি operator
\end{enumerate}

\item Return statement কী করে? \textbf{[ঢাকা-২০২৩]}
\begin{enumerate}[label=\alph*)]
\item Function থেকে মান ফেরত দেয়
\item Function শুরু করে
\item Function বন্ধ করে
\item Function call করে
\end{enumerate}

\item Syntax Error কী? \textbf{[চট্টগ্রাম-২০২২]}
\begin{enumerate}[label=\alph*)]
\item ভাষার নিয়ম ভঙ্গ
\item Logic এ ভুল
\item Runtime এ ভুল
\item কোনো ভুল নয়
\end{enumerate}

\item Logical Error কী? \textbf{[সিলেট-২০২৩]}
\begin{enumerate}[label=\alph*)]
\item Logic এ ভুল
\item ভাষার নিয়ম ভঙ্গ
\item Runtime এ ভুল
\item কোনো ভুল নয়
\end{enumerate}

\item Runtime Error কী? \textbf{[রাজশাহী-২০২২]}
\begin{enumerate}[label=\alph*)]
\item প্রোগ্রাম চলার সময় ভুল
\item ভাষার নিয়ম ভঙ্গ
\item Logic এ ভুল
\item কোনো ভুল নয়
\end{enumerate}

\item Compilation কী? \textbf{[বরিশাল-২০২৩]}
\begin{enumerate}[label=\alph*)]
\item Source code কে machine code এ রূপান্তর
\item Machine code কে source code এ রূপান্তর
\item প্রোগ্রাম চালানো
\item প্রোগ্রাম লেখা
\end{enumerate}

\item Execution কী? \textbf{[দিনাজপুর-২০২২]}
\begin{enumerate}[label=\alph*)]
\item প্রোগ্রাম চালানো
\item প্রোগ্রাম লেখা
\item প্রোগ্রাম compile করা
\item প্রোগ্রাম সংরক্ষণ
\end{enumerate}

\item Library Function কী? \textbf{[কুমিল্লা-২০২৩]}
\begin{enumerate}[label=\alph*)]
\item পূর্বনির্ধারিত function
\item ব্যবহারকারী তৈরি function
\item একটি variable
\item একটি operator
\end{enumerate}

\item User-defined Function কী? \textbf{[যশোর-২০২২]}
\begin{enumerate}[label=\alph*)]
\item ব্যবহারকারী তৈরি function
\item পূর্বনির্ধারিত function
\item একটি variable
\item একটি operator
\end{enumerate}

\item Header File কী? \textbf{[ঢাকা-২০২৩]}
\begin{enumerate}[label=\alph*)]
\item Library function এর declaration
\item প্রোগ্রামের মূল অংশ
\item একটি variable
\item একটি operator
\end{enumerate}

\item Preprocessor Directive কী? \textbf{[চট্টগ্রাম-২০২২]}
\begin{enumerate}[label=\alph*)]
\item \# দিয়ে শুরু হওয়া নির্দেশ
\item একটি function
\item একটি variable
\item একটি operator
\end{enumerate}

\item Macro কী? \textbf{[সিলেট-২০২৩]}
\begin{enumerate}[label=\alph*)]
\item \#define দিয়ে তৈরি constant
\item একটি function
\item একটি variable
\item একটি operator
\end{enumerate}

\item Type Casting কী? \textbf{[রাজশাহী-২০২২]}
\begin{enumerate}[label=\alph*)]
\item এক data type কে অন্যটিতে রূপান্তর
\item ডেটা সংরক্ষণ
\item ডেটা মুছে ফেলা
\item ডেটা প্রসেসিং
\end{enumerate}

\end{enumerate}
\end{multicols}

\subsection*{\textbf{\underline{উত্তরমালা (৫৭-১০৬)}}}

\begin{multicols}{4}
\begin{enumerate}[start=57]
\item a
\item a
\item a
\item a
\item a
\item a
\item a
\item a
\item a
\item a
\item a
\item a
\item a
\item a
\item a
\item a
\item a
\item a
\item a
\item a
\item a
\item a
\item a
\item a
\item a
\item a
\item a
\item a
\item a
\item a
\item a
\item a
\item a
\item a
\item a
\item a
\item a
\item a
\item a
\item a
\item a
\item a
\item a
\item a
\item a
\item a
\item a
\item a
\item a
\item a
\end{enumerate}
\end{multicols}
